\chapter{Tablas de apoyo}

\section{Factores nacionales por mes}

La tabla~\ref{tab:knat} contiene los factores nacionales que sirven de centro de la contracción (ecuación~\eqref{eq:ktilde}). La columna $1/\widehat k^{V}_{\mathrm{nac}}$ dice cuántas detecciones de VIIRS equivalen a un hotspot de MODIS.

\begin{table}[H]
\centering
\caption{Factores nacionales. $\widehat k^{V}_{\mathrm{nac}}$ se calibró con 2012--2019 (sin enero de 2012) y $\widehat k^{T}_{\mathrm{nac}}$, con 2003--2011.}\label{tab:knat}
\small
\begin{tabular}{@{}lrrr@{}}
\toprule
Mes & $\widehat k^{V}_{\mathrm{nac}}$ & $1/\widehat k^{V}_{\mathrm{nac}}$ & $\widehat k^{T}_{\mathrm{nac}}$ \\
\midrule
Enero & 0.292 & 3.42 & 2.574 \\
Febrero & 0.265 & 3.78 & 2.484 \\
Marzo & 0.263 & 3.81 & 2.590 \\
Abril & 0.284 & 3.52 & 2.232 \\
Mayo & 0.247 & 4.04 & 2.125 \\
Junio & 0.195 & 5.14 & 2.145 \\
Julio & 0.221 & 4.52 & 2.176 \\
Agosto & 0.242 & 4.13 & 2.189 \\
Septiembre & 0.302 & 3.31 & 2.066 \\
Octubre & 0.270 & 3.70 & 2.098 \\
Noviembre & 0.251 & 3.98 & 2.340 \\
Diciembre & 0.306 & 3.27 & 2.606

\\
\bottomrule
\end{tabular}
\end{table}

\section{Pronóstico mensual de 2026 para México}

\begin{table}[H]
\centering
\caption{Pronóstico desde el 31 de diciembre de 2025 (hotspots equivalentes a MODIS). El intervalo de $80\,\%$ es de cada mes por separado; el total no lleva intervalo (véase la sección~\ref{sec:interv}). Se incluyen el observado de 2025 y el promedio 2013--2025 como referencia.}\label{tab:pron2026}
\small
\begin{tabular}{@{}lrrrr@{}}
\toprule
Mes & Pronóstico & Intervalo de 80\,\% & Observado 2025 & Promedio 2013--2025 \\
\midrule
Enero 2026 & 2\,298 & 1\,426--3\,382 & 2\,031 & 2\,154 \\
Febrero 2026 & 2\,732 & 1\,697--4\,020 & 2\,346 & 2\,679 \\
Marzo 2026 & 6\,226 & 3\,873--9\,150 & 6\,709 & 6\,011 \\
Abril 2026 & 15\,987 & 9\,956--23\,486 & 13\,755 & 16\,669 \\
Mayo 2026 & 19\,372 & 12\,065--28\,457 & 18\,724 & 20\,844 \\
Junio 2026 & 6\,088 & 3\,788--8\,948 & 3\,302 & 5\,429 \\
Julio 2026 & 1\,346 & 833--1\,984 & 999 & 1\,161 \\
Agosto 2026 & 844 & 520--1\,246 & 797 & 826 \\
Septiembre 2026 & 574 & 352--850 & 427 & 532 \\
Octubre 2026 & 617 & 379--914 & 580 & 618 \\
Noviembre 2026 & 882 & 544--1\,303 & 991 & 793 \\
Diciembre 2026 & 1\,639 & 1\,015--2\,414 & 1\,696 & 1\,632 \\
\midrule
Total & 58\,605 & --- & 52\,357 & 59\,348

\\
\bottomrule
\end{tabular}
\end{table}

\chapter{Preguntas que puede hacer el jurado}

Cada respuesta es corta y señala dónde está el sustento en este documento.

\begin{description}[leftmargin=1em,style=nextline,itemsep=8pt]
  \item[¿Por qué no basta un solo factor para todo México?]
  Porque la razón entre sensores cambia con el mes y con el estado (figura~\ref{fig:knat}). Con un factor único el error por estado y mes es de $24.8\,\%$ y $21.3\,\%$; con factor por estado y mes baja a $12.9\,\%$ y $14.5\,\%$ (tabla~\ref{tab:niveles}).

  \item[¿Cómo sabemos que no está ajustado a los datos que usamos para medir?]
  Medimos con dos ejercicios en los que el año evaluado no entra al cálculo del factor: dejar un año fuera (2012--2019) y predecir 2020--2025, años que el método no usó (capítulo~\ref{cap:valid}). El único parámetro que miró ambos es $\tau$, y su efecto es de décimas de punto (sección~\ref{sec:tau}).

  \item[¿Por qué una razón de totales y no un promedio de razones?]
  Porque es el estimador de máxima verosimilitud condicional bajo un modelo de Poisson y no necesita modelar la intensidad del fuego de cada año (sección~\ref{sec:est-razon}); además da más peso a los años con más datos.

  \item[¿Qué es la contracción y por qué $\tau=100$?]
  Es una media ponderada entre el factor del estado y el nacional del mismo mes, con peso $\sum V/(\sum V+\tau)$ para el estado; equivale a una media a posteriori beta--binomial (sección~\ref{sec:contraccion}). El error es casi igual entre $\tau=30$ y $\tau=300$ (tabla~\ref{tab:tau}).

  \item[¿Por qué descartaron fusionar píxeles de VIIRS y filtrar por FRP?]
  Porque empeoraron todas las métricas (tabla~\ref{tab:variantes}); el error anual fuera de muestra pasó de $2.1\,\%$ a entre $4.4\,\%$ y $7.6\,\%$.

  \item[¿Qué tan confiable es la parte de 2000--2002?]
  Es la más incierta: el factor de Terra se calibró con 2003--2011 y se aplica hacia atrás sin posibilidad de comprobarlo (sección~\ref{sec:terra}).

  \item[¿El pronóstico mejora a la climatología?]
  A escala nacional no de manera apreciable (tabla~\ref{tab:pron-nac}); sí supera con claridad a repetir el año anterior. Por estado supera a la climatología en 22 a 26 de los 32 estados según el horizonte, con diferencias medianas pequeñas (tabla~\ref{tab:pron-est}).

  \item[¿El pronóstico de 2026 es correcto?]
  No se puede saber con nuestros datos, que terminan en diciembre de 2025. Lo que sí se midió es el desempeño del método de 2020 a 2025 (capítulo~\ref{cap:pron}).

  \item[¿Por qué los intervalos cubren más de $80\,\%$?]
  A escala nacional cubren de $85$ a $89\,\%$: son algo conservadores. Por estado la cobertura mediana es de $82$ a $85\,\%$, más cercana a lo nominal.

  \item[¿Un hotspot es una hectárea quemada?]
  No. Es un píxel con anomalía térmica; no mide superficie quemada ni daño (capítulo~\ref{cap:limites}).
\end{description}

\chapter{Glosario}

\begin{description}[leftmargin=1em,style=nextline,itemsep=6pt]
  \item[Armonizar.] Llevar los conteos de dos sensores a una unidad común de modo que sean comparables.
  \item[Climatología.] Promedio histórico de una variable para la misma fecha del calendario.
  \item[Contracción.] Acercar una estimación ruidosa hacia un valor de referencia más estable, con un peso que depende de la cantidad de datos.
  \item[FIRMS.] Sistema de la NASA que distribuye detecciones de fuego activo casi en tiempo real.
  \item[FRP.] Potencia radiativa del fuego, en megavatios.
  \item[Hotspot.] Píxel de satélite con anomalía térmica compatible con fuego activo.
  \item[Hotspot equivalente a MODIS.] Unidad común de la serie armonizada: lo que habrían contado Terra y Aqua.
  \item[LOYO.] Validación cruzada que deja un año fuera (\emph{leave-one-year-out}).
  \item[Origen.] Último día observado a partir del cual se pronostica.
  \item[Sesgo.] Tendencia sistemática de un método a sobrestimar o subestimar.
  \item[Sobredispersión.] Variabilidad de los conteos mayor que la de un modelo de Poisson.
  \item[WAPE.] Error porcentual absoluto ponderado: suma de errores absolutos entre suma de lo observado.
\end{description}
